\scanpage{73}\mspage{72}
\uline{Cor 4}\quad \inserted{Soit $F$ une figure\unsure{,|.}} Pour qu'une partie $A$ de
$\unsure{\mathcal L|L}$ \unsure{définisse} une figure discrète\inserted{$\unsure{\Phi|\varphi}$} et qu'on
ait $\Phi\ll F$, il \unsure{f. et. s.|f. et s.} que
\[
\begin{aligned}
\text{(i)}\quad &\Phi\subset\mathrm{omb}(F)\\
\text{(ii)}\quad &\forall X\in\widetilde F,\qquad \Phi\cap\mathrm{omb}(X)\quad \text{``soit une figure''}
\end{aligned}
\]
(p.ex soit fini et \unsure{formé} de \deleted{\unsure{élém.|él}}\deleted{\unsure{de|d'}} \unsure{lieux} deux à deux
disjoints).

Ceci implique \unsure{notamment} que à la
page \unsure{67|62}, dans l'énoncé de \unsure{At|Ax} $9'$
(critère suffisant pour que $X\mathrel{\rlap{$<$}{>}}Y$), la
condition donnée \deleted{soit} en termes
de disjonction de lieux, est aussi
\uline{nécessaire} [mais cela résulte aussi directement
de At 10]

% slanted
\margin{At 10 \unsure{posé} après, \unsure{page} 64.}
% two left pen rules: the outer one spans statement and proof
\noindent\llap{\smash{\rule[-9.2\baselineskip]{0.4pt}{9.8\baselineskip}}\quad}%
\uline{Corollaire $\deleted{\unsure{4}}5$}\quad Soient $F$ une figure, $G$ une
\llap{\smash{\rule[-1.4\baselineskip]{0.4pt}{1.8\baselineskip}}\enspace}sous figure $\neq F$. Alors $\exists x\in\mathcal L$ tel que
$x\unsure{\ll|\underline{\ll}}\unsure{\mathbb F|F}$ (i.e. $x\in\mathrm{omb}(\mathbb F)$) et $x\parallel G$.

Soit $X\in\widetilde{\mathbb F}\setminus\widetilde G$, et $x\in\mathrm{omb}(X)^\circ$ \unsure{(148)|(At 8)|(A48)|(1.48)}
\unsure{i.e.|i.e.e} $x\in\mathcal L$, $x\unsure{\overset{\circ}{\ll}|\overset{\circ}{\underline{\ll}}}X$. Je dis qu'il fait l'affaire.
Pour voir qu'il est disjoint de $G$, on peut
soit invoquer At 9 (critère de disjonction pour les
lieux) et At 10, (critère de disjonction pour les lieux),
soit se \unsure{référer} à \deleted{\unsure{un|ce} \unsure{même|axiome}} l'axiome de recollement, pour
\unsure{construire} $G\amalg x$ comme \unsure{raffinement de|raffinant} $F$ ...
